% ==============================================================================
% CHƯƠNG 1: TỔNG QUAN VÀ TÌM HIỂU CÁC CÔNG NGHỆ SỬ DỤNG
% ==============================================================================
\chapter{TỔNG QUAN VÀ TÌM HIỂU CÁC CÔNG NGHỆ SỬ DỤNG}
\label{chap:tong_quan_cong_nghe}

\section{Tổng quan về đề tài}
\label{sec:tong_quan_de_tai}

\subsection{Khái niệm hệ thống POS và vai trò trong bán lẻ hiện đại}
Hệ thống Điểm bán lẻ (Point of Sale - viết tắt là POS) là tập hợp các giải pháp phần cứng và phần mềm hỗ trợ doanh nghiệp và các cửa hàng bán lẻ thực hiện các giao dịch thương mại thanh toán hàng hóa, dịch vụ trực tiếp tại quầy thu ngân. Trong lịch sử phát triển của ngành bán lẻ, các giao dịch ban đầu phụ thuộc hoàn toàn vào sổ sách kế toán thủ công hoặc các máy tính tiền cơ học (Mechanical Cash Registers). Các phương pháp truyền thống này bộc lộ nhiều điểm hạn chế nghiêm trọng: tốc độ thanh toán chậm, dễ xảy ra sai sót khi tính tiền và trả lại tiền thừa, không có khả năng kiểm soát số lượng hàng hóa tồn kho theo thời gian thực, và đặc biệt là không thể ngăn chặn các hành vi gian lận nội bộ.

Trong kỷ nguyên chuyển đổi số và thương mại hiện đại, hệ thống POS đã tiến hóa trở thành trung tâm điều phối dữ liệu (Data Hub) của toàn bộ hoạt động kinh doanh tại cửa hàng. Một hệ thống POS hoàn chỉnh không đơn thuần chỉ đóng vai trò là công cụ tính tổng tiền đơn hàng, mà còn tích hợp các phân hệ nghiệp vụ cốt lõi:
\begin{itemize}[leftmargin=1.5cm]
    \item \textbf{Tự động hóa thanh toán tại quầy:} Hỗ trợ quét mã vạch sản phẩm, tính toán tự động các loại thuế (VAT), chiết khấu giảm giá, chương trình khuyến mãi, và hỗ trợ đa dạng phương thức thanh toán từ tiền mặt truyền thống đến chuyển khoản ngân hàng qua mã QR động (VietQR).
    \item \textbf{Quản lý tồn kho theo thời gian thực (Real-time Inventory Tracking):} Cập nhật tức thời số dư tồn kho ngay khi giao dịch bán lẻ hoàn tất thông qua cơ chế sổ cái kế toán (Ledger-based Inventory), ngăn chặn tình trạng thất thoát hoặc bán vượt quá số lượng hàng có sẵn trong kho.
    \item \textbf{Quản trị quan hệ khách hàng (CRM):} Nhận diện khách hàng thân thiết thông qua số điện thoại, tự động tích lũy lịch sử chi tiêu, số lần mua hàng và điểm thưởng nhằm phục vụ các chiến dịch chăm sóc khách hàng cá nhân hóa.
    \item \textbf{Báo cáo và phân tích kinh doanh (Business Intelligence):} Cung cấp các chỉ số đo lường hiệu quả hoạt động (KPI) trực quan như doanh thu theo ca làm việc, biểu đồ doanh thu theo ngày/tháng, và danh mục các mặt hàng bán chạy nhất để người quản lý ra quyết định nhập hàng kịp thời.
\end{itemize}

\subsection{Sự cần thiết của việc xây dựng hệ thống POS trên nền tảng Web}
Trước đây, đa số các phần mềm POS được xây dựng dưới dạng ứng dụng cài đặt cục bộ trên máy tính để bàn (Desktop Application) chạy trên hệ điều hành Windows. Mặc dù ứng dụng Desktop có ưu thế về tốc độ phản hồi đồ họa ngoại tuyến, song mô hình này bộc lộ những rào cản kỹ thuật rất lớn trong bối cảnh chuỗi cửa hàng và thương mại đa kênh:
\begin{enumerate}[label=\arabic*., leftmargin=1.5cm]
    \item \textbf{Phụ thuộc chặt chẽ vào phần cứng:} Ứng dụng Desktop đòi hỏi máy tính thu ngân phải có cấu hình tương thích, việc cài đặt và nâng cấp phần mềm phải thực hiện thủ công trên từng máy trạm, làm gia tăng chi phí vận hành và bảo trì kỹ thuật.
    \item \textbf{Dữ liệu bị cô lập (Data Silo):} Khi cơ sở dữ liệu được lưu cục bộ tại quầy, chủ cửa hàng không thể theo dõi tình hình kinh doanh từ xa trừ khi thiết lập các kết nối mạng riêng ảo (VPN) phức tạp và kém ổn định.
    \item \textbf{Khó khăn trong việc tích hợp dịch vụ bên thứ ba:} Việc tích hợp các cổng thanh toán điện tử thế hệ mới, dịch vụ lưu trữ đám mây hoặc hệ thống phát hành hóa đơn điện tử đòi hỏi cấu hình mạng phức tạp và tiềm ẩn nhiều rủi ro an ninh mạng.
\end{enumerate}

Để khắc phục triệt để các hạn chế trên, việc xây dựng hệ thống POS trên kiến trúc Web phân tán (Web-based POS Architecture) là xu thế tất yếu. Ứng dụng Web POS cho phép nhân viên thu ngân truy cập và thực hiện thao tác bán hàng thông qua bất kỳ trình duyệt tiêu chuẩn nào (Google Chrome, Microsoft Edge, Safari) trên nhiều thiết bị phần cứng khác nhau (máy tính để bàn, máy tính bảng, màn hình cảm ứng POS chuyên dụng). Đồng thời, kiến trúc Web tách biệt giữa tầng Giao diện (Front-end) và tầng Dịch vụ máy chủ (Back-end) thông qua các giao diện lập trình ứng dụng RESTful API giúp hệ thống dễ dàng mở rộng, đồng bộ dữ liệu tức thời lên máy chủ trung tâm và cho phép người quản lý kiểm soát toàn bộ hoạt động kinh doanh từ xa ở mọi thời điểm \cite{sommerville2016}.

\subsection{Mục tiêu và phạm vi của đề tài}
Đề tài \textbf{''Xây dựng website POS Quản lý bán hàng''} thuộc khuôn khổ môn học \textit{Phát triển phần mềm mã nguồn mở} tại Trường Đại học Giao thông Vận tải hướng tới các mục tiêu cụ thể:
\begin{itemize}[leftmargin=1.5cm]
    \item Nghiên cứu và ứng dụng toàn diện các công nghệ và thư viện mã nguồn mở hiện đại hàng đầu hiện nay, bao gồm Java Spring Boot 4.x, ReactJS 19, Vite, Tailwind CSS v4, MySQL Server và Nginx.
    \item Thiết kế và cài đặt một hệ thống POS hoàn chỉnh, đáp ứng tiêu chuẩn khắt khe về thời gian phản hồi tại quầy thu ngân (dưới $1$ giây cho mỗi thao tác quét và thêm mặt hàng vào giỏ).
    \item Xây dựng cơ chế thanh toán không dùng tiền mặt thế hệ mới thông qua tích hợp cổng thanh toán VietQR động của PayOS, kết hợp quy trình xử lý giao dịch bù trừ (Compensating Transaction) đảm bảo tính toàn vẹn số liệu tài chính và tồn kho kế toán.
    \item Đóng gói và triển khai ứng dụng bằng công nghệ containerization Docker, đảm bảo khả năng tái lập môi trường và tính khả chuyển cao trên mọi hạ tầng máy chủ đám mây.
\end{itemize}

% ==============================================================================
\section{Công nghệ giao diện người dùng (Front-end)}
\label{sec:cong_nghe_frontend}

\subsection{Nền tảng ReactJS (Phiên bản React 19)}
ReactJS là một thư viện mã nguồn mở bằng ngôn ngữ JavaScript do Meta (Facebook) phát triển và duy trì, chuyên dùng để xây dựng các giao diện người dùng tương tác cao (Interactive User Interfaces) dựa trên kiến trúc hướng thành phần (Component-Based Architecture) \cite{banks2020}. 

Trong hệ thống POS, việc lựa chọn React 19 mang lại những ưu thế vượt trội:
\begin{itemize}[leftmargin=1.5cm]
    \item \textbf{Virtual DOM (Mô hình cây tài liệu ảo):} Khi trạng thái giỏ hàng hoặc danh mục sản phẩm thay đổi, React không cập nhật trực tiếp toàn bộ cây DOM thực của trình duyệt (vốn là thao tác tốn kém tài nguyên tính toán). Thay vào đó, thuật toán Reconciliation (Fiber Engine) tiến hành so sánh đối chiếu (diffing) giữa Virtual DOM mới và cũ để chỉ cập nhật chính xác các phần tử HTML bị thay đổi. Điều này giúp giao diện POS phản hồi mượt mà ngay cả khi giỏ hàng có hàng chục mặt hàng và liên tục thay đổi số lượng.
    \item \textbf{React Hooks hiện đại:} Việc sử dụng các hook chuẩn hóa như \texttt{useState}, \texttt{useReducer}, \texttt{useCallback} và \texttt{useMemo} giúp đóng gói logic giao diện gọn gàng, loại bỏ hoàn toàn sự phức tạp của các Class Component truyền thống, tối ưu hóa quá trình tính toán lại tổng tiền và chiết khấu.
    \item \textbf{Cấu trúc Component hóa (Modular Components):} Cho phép chia nhỏ giao diện bán hàng thành các khối độc lập có khả năng tái sử dụng cao như: \texttt{DisplayCategories}, \texttt{DisplayItems}, \texttt{CartItems}, \texttt{POSItemModal}, \texttt{ReceiptPopup}.
\end{itemize}

\subsection{Công cụ đóng gói và máy chủ phát triển Vite 8}
Trong quá trình phát triển các ứng dụng Front-end hiện đại, việc sử dụng các bộ đóng gói mã nguồn (Bundlers) truyền thống như Webpack thường gặp phải hiện tượng nghẽn cổ chai: thời gian khởi động máy chủ thử nghiệm (Cold Start) rất chậm và thời gian cập nhật mô-đun nóng (Hot Module Replacement - HMR) bị trễ đáng kể khi quy mô dự án phình to.

Dự án lựa chọn \textbf{Vite 8} làm nền tảng công cụ phát triển. Vite tận dụng triệt để tính năng \textit{Native ES Modules} (ESM) sẵn có trên các trình duyệt hiện đại. Thay vì phải đóng gói toàn bộ mã nguồn trước khi khởi động, Vite chỉ biên dịch các mô-đun mã nguồn theo yêu cầu (On-demand Compilation) bằng công cụ esbuild viết bằng ngôn ngữ Go. Khi mã nguồn của một component trong giao diện POS được chỉnh sửa, Vite chỉ gửi tín hiệu cập nhật chính xác file đó qua WebSocket, giúp quá trình phản chiếu thay đổi lên giao diện diễn ra gần như tức thì (dưới 50 mili-giây) \cite{vitedoc2024}. Khi đóng gói cho môi trường sản xuất (Production Build), Vite sử dụng Rollup để tạo ra các gói tệp tĩnh (bundle) được nén tối ưu, phân chia mã nguồn thông minh (Code Splitting) và loại bỏ hoàn toàn các đoạn mã thừa không sử dụng (Tree-shaking).

\subsection{Framework định kiểu Tailwind CSS (Phiên bản v4)}
Thay vì sử dụng các bộ UI Component được tạo sẵn với kiểu dáng cố định (như Bootstrap hay Ant Design) vốn chứa nhiều lớp CSS thừa làm tăng kích thước tệp tải về, dự án áp dụng \textbf{Tailwind CSS v4} kết hợp cùng plugin chính thức \texttt{@tailwindcss/vite}.

Tailwind CSS hoạt động theo triết lý \textit{Utility-First CSS}, cung cấp hàng nghìn lớp tiện ích cấp thấp (Utility Classes) có thể kết hợp linh hoạt trực tiếp trong mã JSX:
\begin{itemize}[leftmargin=1.5cm]
    \item \textbf{Tối ưu hiệu năng nạp trang:} Engine xử lý của Tailwind CSS v4 phân tích cú pháp tĩnh các tệp mã nguồn để chỉ sinh ra các lớp CSS thực sự được sử dụng. Nhờ đó, kích thước file CSS khi xuất xưởng thường nhỏ hơn $20$ KB, loại bỏ hoàn toàn gánh nặng băng thông tải trang.
    \item \textbf{Tính nhất quán về Design Tokens:} Dự án xây dựng bảng màu nhận diện thương hiệu chuyên nghiệp cho POS tại tệp \texttt{src/index.css}: Màu chủ đạo Primary (\#2563EB - Blue 600), màu điểm nhấn Accent (\#059669 - Emerald 600) biểu thị giao dịch thành công, màu cảnh báo Danger (\#DC2626 - Red 600) cho thao tác hủy đơn, và phông chữ Google Font \textit{Outfit} mang phong cách hiện đại.
    \item \textbf{Hỗ trợ định dạng in nhiệt chuyên biệt:} Tailwind cho phép định nghĩa các lớp dành riêng cho môi trường in ấn (\texttt{@media print}), giúp mẫu hóa đơn thanh toán \texttt{ReceiptPopup} tự động ẩn đi toàn bộ khung viền, thanh điều hướng và chỉ hiển thị đúng nội dung phiếu in 80mm tương thích với các máy in nhiệt chuyên dụng.
\end{itemize}

\subsection{Quản lý trạng thái máy chủ với TanStack React Query (v5)}
Một trong những thách thức kỹ thuật lớn nhất của ứng dụng POS là sự khác biệt giữa hai loại trạng thái dữ liệu:
\begin{enumerate}[label=\alph*., leftmargin=1.5cm]
    \item \textbf{Client State (Trạng thái cục bộ tại máy khách):} Gồm giỏ hàng hiện tại, trạng thái đóng/mở của các cửa sổ modal, giá trị nhập trong ô tìm kiếm. Trạng thái này được quản lý thông qua các React Hook (\texttt{useState}, \texttt{useCartItem}).
    \item \textbf{Server State (Trạng thái dữ liệu phía máy chủ):} Gồm danh mục hàng hóa, số dư tồn kho của các biến thể (SKU), danh sách chương trình khuyến mãi đang kích hoạt, và đặc biệt là trạng thái thanh toán của đơn hàng chuyển khoản. Trạng thái này thuộc quyền sở hữu của máy chủ và có thể bị thay đổi bất ngờ bởi các tác nhân khác ngoài tầm kiểm soát của trình duyệt.
\end{enumerate}

Để giải quyết bài toán này mà không phải dựa vào các thư viện quản lý trạng thái cồng kềnh như Redux hay MobX, dự án áp dụng \textbf{TanStack React Query v5} \cite{tanstackquery2024}. React Query đảm nhận toàn bộ vòng đời của dữ liệu máy chủ:
\begin{itemize}[leftmargin=1.5cm]
    \item \textbf{Cơ chế bộ nhớ đệm thông minh (Caching \& Stale-While-Revalidate):} Danh mục sản phẩm và biến thể được lưu trong bộ nhớ tạm client. Khi nhân viên điều hướng giữa các màn hình, dữ liệu hiển thị ngay lập tức từ bộ nhớ đệm trong khi request nền vẫn được kích hoạt để kiểm tra tính cập nhật.
    \item \textbf{Tự động vô hiệu hóa bộ nhớ đệm (Cache Invalidation):} Khi thực hiện một thao tác ghi (Mutation) như thêm mới sản phẩm, nhập kho hoặc hủy hóa đơn, React Query tự động đánh dấu các khóa truy vấn liên quan (\texttt{['items']}, \texttt{['orders']}) là hết hạn và âm thầm kích hoạt tải lại dữ liệu mới nhất.
    \item \textbf{Thăm dò tự động (Polling Mechanism):} Trong luồng thanh toán chuyển khoản VietQR PayOS, hook tùy biến \texttt{usePolledOrderData} thiết lập cơ chế gửi truy vấn thăm dò với chu kỳ $3$ giây/lần. Khi khách hàng quét mã và chuyển tiền thành công, trạng thái đơn hàng trên màn hình thu ngân tự động chuyển sang \texttt{COMPLETED} mà không yêu cầu nhân viên phải nhấn nút F5 tải lại trang.
\end{itemize}

\subsection{Xử lý biểu mẫu với React Hook Form}
Tại các phân hệ quản trị (thêm mới sản phẩm có nhiều biến thể SKU, cấu hình nhóm tùy chọn Modifier, thiết lập chương trình khuyến mãi BOGO), số lượng trường dữ liệu cần nhập trên một biểu mẫu là rất lớn. Việc sử dụng phương pháp quản lý biểu mẫu truyền thống (Controlled Components - mỗi lần gõ phím cập nhật một state) sẽ khiến toàn bộ Form bị Re-render liên tục, gây ra hiện tượng giật lag rõ rệt trên các máy POS cấu hình yếu.

Dự án tích hợp \textbf{React Hook Form v7}, áp dụng kỹ thuật \textit{Uncontrolled Components} thông qua việc gán trực tiếp tham chiếu DOM (\texttt{ref}). React Hook Form chỉ đọc giá trị của các trường khi người dùng thực hiện thao tác gửi (Submit) hoặc khi cần kiểm tra tính hợp lệ (Validation). Kết hợp cùng tính năng \texttt{useFieldArray}, hệ thống cho phép người dùng thêm/bớt động danh sách các biến thể hàng hóa và thuộc tính JSON một cách mượt mà và tiêu tốn tối thiểu tài nguyên bộ nhớ.

\subsection{So sánh và tổng kết giải pháp Front-end}
Bảng \ref{tab:so_sanh_frontend} tổng hợp so sánh các giải pháp công nghệ Front-end phổ biến, làm rõ cơ sở khoa học của việc lựa chọn React 19 + Vite + Tailwind CSS + React Query cho đề tài.

\begin{table}[htbp]
\centering
\caption{So sánh các giải pháp công nghệ Front-end cho hệ thống POS}
\label{tab:so_sanh_frontend}
\vspace{2mm}
\begin{tabularx}{\textwidth}{|p{3cm}|X|X|X|}
\hline
\textbf{Tiêu chí so sánh} & \textbf{ReactJS + Vite + Tailwind} & \textbf{VueJS + Nuxt} & \textbf{Angular Framework} \\ \hline
\textbf{Kiến trúc cốt lõi} & Thư viện linh hoạt, Virtual DOM tối ưu, đóng gói Native ESM siêu tốc. & Khung làm việc tiến bộ, hệ thống phản ứng (Reactivity System) 2 chiều. & Khung làm việc hoàn chỉnh (All-in-one), phụ thuộc Dependency Injection nặng. \\ \hline
\textbf{Tốc độ phản hồi quầy} & \textbf{Rất cao} ($< 50$ms), chi phí re-render cực thấp nhờ React Query và Hook Form. & Cao, giao diện mượt mà nhưng hệ sinh thái thư viện ngoài phân mảnh hơn. & Trung bình, kích thước bundle ban đầu lớn làm chậm thời gian tải trang đầu tiên. \\ \hline
\textbf{Khả năng tùy biến giao diện POS} & \textbf{Tuyệt đối}, Tailwind CSS cho phép can thiệp sâu vào từng điểm ảnh và mẫu in nhiệt 80mm. & Tốt, nhưng thường bị ràng buộc bởi các component scoped CSS có sẵn. & Kém linh hoạt hơn, tốn nhiều công sức ghi đè các khuôn mẫu CSS của Angular Material. \\ \hline
\textbf{Cộng đồng \& Tài liệu} & Hệ sinh thái mã nguồn mở lớn nhất thế giới, cộng đồng hỗ trợ đông đảo. & Cộng đồng lớn mạnh, tài liệu rõ ràng nhưng ít giải pháp chuyên sâu cho POS. & Hướng tới doanh nghiệp lớn truyền thống, lộ trình học tập dốc. \\ \hline
\end{tabularx}
\end{table}

% ==============================================================================
\section{Công nghệ xử lý nghiệp vụ máy chủ (Back-end)}
\label{sec:cong_nghe_backend}

\subsection{Ngôn ngữ lập trình Java (Phiên bản Java 25)}
Java là một trong những ngôn ngữ lập trình hướng đối tượng mạnh mẽ, an toàn và phổ biến nhất trên thế giới dành cho các hệ thống phần mềm cấp doanh nghiệp (Enterprise Systems). Dự án áp dụng phiên bản **Java 25** (bản phát hành mới nhất thuộc chu kỳ phát triển mã nguồn mở OpenJDK) với hàng loạt cải tiến quan trọng:
\begin{itemize}[leftmargin=1.5cm]
    \item \textbf{Cơ chế quản lý bộ nhớ và thu gom rác tối ưu (Z Garbage Collector):} Đảm bảo thời gian dừng ứng dụng (Stop-The-World Pause Time) dưới $1$ mili-giây ngay cả khi hệ thống xử lý khối lượng lớn các giao dịch ghi log và thanh toán đồng thời.
    \item \textbf{Cú pháp ngôn ngữ hiện đại (Modern Language Features):} Tận dụng các tính năng Record Classes, Pattern Matching, Switch Expressions và khối văn bản đa dòng (Text Blocks) giúp mã nguồn xử lý nghiệp vụ tại tầng Service trở nên ngắn gọn, dễ bảo trì và hạn chế tối đa các lỗi ngoại lệ con trỏ rỗng (\texttt{NullPointerException}).
    \item \textbf{Hệ thống kiểu dữ liệu tĩnh mạnh (Strong Static Typing):} Đảm bảo tính toán tài chính chính xác tuyệt đối thông qua kiểu số thực dấu phẩy động độ chính xác tùy biến \texttt{BigDecimal}, loại trừ hoàn toàn các sai số làm tròn số tiền thường gặp ở các ngôn ngữ có kiểu dữ liệu động như JavaScript hay Python.
\end{itemize}

\subsection{Framework Spring Boot 4.1.0 và hệ sinh thái Spring}
Spring Boot là một framework mã nguồn mở xây dựng dựa trên nền tảng Spring Framework, được thiết kế nhằm đơn giản hóa tối đa quy trình khởi tạo, cấu hình và phát triển các dịch vụ Web độc lập (Production-ready Spring Applications) \cite{walls2022}. 

Các thành phần cốt lõi của Spring Boot được khai thác triệt để trong dự án:
\begin{itemize}[leftmargin=1.5cm]
    \item \textbf{Đảo ngược điều khiển (IoC) và Tiêm phụ thuộc (Dependency Injection - DI):} Mọi thành phần Controller, Service, Repository đều được Spring Container quản lý vòng đời dưới dạng các Bean (\texttt{@Service}, \texttt{@RestController}, \texttt{@Repository}). Cơ chế này giúp các mô-đun nghiệp vụ giảm thiểu tối đa sự phụ thuộc chặt chẽ (Loose Coupling), tạo điều kiện thuận lợi cho việc viết các ca kiểm thử tự động (Unit Test / Integration Test).
    \item \textbf{Spring Data JPA:} Tự động sinh ra mã nguồn thực thi các thao tác cơ sở dữ liệu dựa trên các quy ước đặt tên phương thức trong Interface (Method Name Queries), loại bỏ hàng nghìn dòng mã SQL lặp lại thủ công.
    \item \textbf{Máy chủ nhúng (Embedded Servlet Container):} Ứng dụng tích hợp sẵn máy chủ Apache Tomcat nhúng, cho phép đóng gói toàn bộ dịch vụ backend thành một tệp JAR duy nhất có khả năng thực thi độc lập mà không yêu cầu cài đặt máy chủ ứng dụng bên ngoài.
\end{itemize}

\subsection{Kiến trúc RESTful API đa tầng}
Hệ thống Back-end tuân thủ nghiêm ngặt mô hình kiến trúc phân tầng (Multi-tier Architecture) nhằm đảm bảo nguyên lý đơn trách nhiệm (Single Responsibility Principle) \cite{fowler2003}, như minh họa tại Hình \ref{fig:backend_layers}:

\begin{figure}[htbp]
\centering
\begin{minipage}{0.85\textwidth}
\centering
\fbox{
    \parbox{0.95\textwidth}{
        \centering \vspace{2mm}
        \textbf{MÔ HÌNH KIẾN TRÚC PHÂN TẦNG BACK-END (SPRING BOOT)}\\[3mm]
        \textbf{[Client Browser / Mobile / PayOS Webhook]}\\[1mm]
        $\Downarrow$ \textit{HTTP/JSON Request qua cổng /api/v1.0}\\[1mm]
        \colorbox{blue!10}{\parbox{0.9\textwidth}{\centering \textbf{TẦNG ĐIỀU KHIỂN (REST CONTROLLER LAYER)}\\
        \textit{Xác thực quyền hạn RBAC, tiếp nhận JSON Payload, điều hướng nghiệp vụ}}}\\[1mm]
        $\Downarrow$ \textit{Truyền dữ liệu qua các DTO (Data Transfer Objects)}\\[1mm]
        \colorbox{green!10}{\parbox{0.9\textwidth}{\centering \textbf{TẦNG NGHIỆP VỤ (SERVICE LAYER)}\\
        \textit{Tính toán hóa đơn, trừ kho sổ cái, đánh giá khuyến mãi, bù trừ hoàn đơn}}}\\[1mm]
        $\Downarrow$ \textit{Thao tác nghiệp vụ thực thể}\\[1mm]
        \colorbox{orange!10}{\parbox{0.9\textwidth}{\centering \textbf{TẦNG TRUY XUẤT DỮ LIỆU (REPOSITORY LAYER - JPA)}\\
        \textit{Giao tiếp CSDL qua Spring Data JPA Interfaces và Hibernate ORM}}}\\[1mm]
        $\Downarrow$ \textit{SQL Queries có giao dịch ACID}\\[1mm]
        \colorbox{gray!20}{\parbox{0.9\textwidth}{\centering \textbf{HỆ QUẢN TRỊ CSDL RELATIONAL (MYSQL 8.0)}}}\\[2mm]
    }
}
\end{minipage}
\caption{Kiến trúc phân tầng của hệ thống Back-end Spring Boot}
\label{fig:backend_layers}
\end{figure}

Chức năng chi tiết của từng tầng:
\begin{enumerate}[label=\arabic*., leftmargin=1.5cm]
    \item \textbf{Tầng Trình diễn / Điều khiển (Controller):} Gồm 13 REST Controllers (như \texttt{OrderController}, \texttt{ItemController}, \texttt{PromotionPublicController},...). Tầng này chịu trách nhiệm định tuyến các yêu cầu HTTP theo đúng tiền tố Context Path (\texttt{/api/v1.0}), kiểm tra tính hợp lệ của dữ liệu đầu vào (\texttt{@Valid}) và chuyển giao xử lý cho tầng Service. Theo chuẩn thiết kế của dự án, tầng Controller trả về trực tiếp các đối tượng dữ liệu thô (Raw DTO) hoặc mảng danh sách JSON, không sử dụng các lớp vỏ bọc phức tạp (Wrapper Envelope), giúp việc xử lý dữ liệu tại Front-end diễn ra nhanh chóng và minh bạch.
    \item \textbf{Tầng Nghiệp vụ cốt lõi (Service):} Gồm 11 lớp hiện thực giao diện (Implementations) đảm nhiệm toàn bộ quy tắc logic của bài toán POS:
    \begin{itemize}
        \item \texttt{OrderServiceImpl}: Tái thẩm định giá bán sản phẩm, tính toán lại tổng tiền và chiết khấu phía server nhằm chống gian lận thay đổi giá từ client, ghi sổ cái trừ kho tức thời và khởi tạo link thanh toán PayOS.
        \item \texttt{PromotionServiceImpl}: Thực thi thuật toán đánh giá chương trình khuyến mãi (Promotion Engine) lựa chọn duy nhất một ưu đãi có giá trị giảm cao nhất cho khách hàng theo nguyên tắc không cộng dồn (No Stacking - ADR-0003).
        \item \texttt{InventoryServiceImpl}: Thực hiện quản lý biến động kho theo nguyên tắc kế toán kép (Ledger Transactions).
        \item \texttt{ActivityLogServiceImpl}: Tự động trích xuất danh tính người dùng từ ngữ cảnh an ninh (\texttt{SecurityContextHolder}) để ghi nhận dấu vết kiểm toán (Audit Trail) cho toàn bộ các thao tác thêm, sửa, xóa dữ liệu nhạy cảm.
    \end{itemize}
    \item \textbf{Tầng Truy xuất dữ liệu (Repository):} Kế thừa từ \texttt{JpaRepository<T, ID>}, cung cấp các phương thức truy vấn tối ưu được thực hiện tự động trong các giao dịch có quản lý (\texttt{@Transactional}).
    \item \textbf{Tầng Mô hình dữ liệu và Đối tượng truyền tải (Entity \& DTO):}
    \begin{itemize}
        \item \textit{JPA Entities:} Đại diện cho cấu trúc bảng vật lý trong cơ sở dữ liệu MySQL, được quản lý vòng đời bởi Hibernate EntityManager.
        \item \textit{DTOs (Data Transfer Objects):} Phân tách rạch ròi giữa dữ liệu lưu trữ nội bộ và dữ liệu trao đổi qua mạng (như \texttt{ItemResponse}, \texttt{OrderResponse}, \texttt{AuthRequest}), bảo vệ an toàn các thông tin nhạy cảm (như mật khẩu đã băm hoặc token nội bộ).
    \end{itemize}
\end{enumerate}

\subsection{Cơ chế bảo mật, xác thực và phân quyền (Spring Security)}
Bảo mật là một trong những thành tố quan trọng nhất của một hệ thống quản lý bán hàng có phân quyền nhân viên và xử lý tiền tệ. Dự án triển khai hệ thống an ninh đa lớp dựa trên \textbf{Spring Security 6.x} kết hợp chuẩn \textbf{JWT (JSON Web Token)} \cite{rfc7519} và quyết định kiến trúc \textbf{ADR-0007}:

\subsubsection{Phân quyền người dùng dựa trên vai trò (RBAC)}
Hệ thống thiết lập 2 vai trò quyền hạn riêng biệt:
\begin{itemize}[leftmargin=1.5cm]
    \item \textbf{Vai trò Quản trị viên (\texttt{ROLE\_ADMIN}):} Có toàn quyền truy cập các API quản trị bắt đầu bằng tiền tố \texttt{/api/v1.0/admin/**}, bao gồm: đăng ký/sửa/xóa tài khoản nhân viên, điều chỉnh danh mục hàng hóa, tạo mã khuyến mãi, nhập xuất kho thủ công và kiểm kê cân bằng kho.
    \item \textbf{Vai trò Nhân viên thu ngân (\texttt{ROLE\_STAFF}):} Chỉ được cấp quyền thực hiện các thao tác phục vụ trực tiếp việc bán hàng tại quầy: xem sản phẩm, tạo hóa đơn bán lẻ, đổi phương thức thanh toán sang tiền mặt, in hóa đơn, tra cứu khách hàng và chỉ được xem nhật ký hoạt động do chính tài khoản của mình tạo ra.
\end{itemize}

\subsubsection{Kiến trúc bảo mật hai lớp Token chống Replay Attack (ADR-0007)}
Nhằm triệt tiêu hoàn toàn nguy cơ bị đánh cắp phiên đăng nhập qua các cuộc tấn công Cross-Site Scripting (XSS), hệ thống không lưu trữ bất kỳ token xác thực nào vào \texttt{localStorage} hay \texttt{sessionStorage} của trình duyệt:
\begin{enumerate}[label=\arabic*., leftmargin=1.5cm]
    \item \textbf{Access Token (Khóa truy cập ngắn hạn):} Có thời hạn hiệu lực là $15$ phút ($900.000$ ms). Token này chỉ được lưu trữ tạm thời trong bộ nhớ RAM của ứng dụng Front-end (\texttt{authSession.js}) và được đính kèm vào tiêu đề HTTP \texttt{Authorization: Bearer <token>} thông qua bộ can thiệp Axios Interceptor trên mỗi request.
    \item \textbf{Refresh Token (Khóa làm mới dài hạn):} Có thời hạn $7$ ngày ($604.800.000$ ms). Token được sinh ngẫu nhiên an toàn (48 bytes cryptographically secure random), băm bằng thuật toán SHA-256 trước khi lưu vào bảng \texttt{refresh\_tokens} trong cơ sở dữ liệu. Token thô được gửi về client duy nhất qua một **Cookie có cờ bảo mật HttpOnly** (\texttt{refreshToken}), với các thuộc tính bảo vệ nghiêm ngặt: \texttt{SameSite=Lax}, \texttt{Path=/api/v1.0}, ngăn ngừa tuyệt đối việc mã JavaScript độc hại trên trình duyệt đọc được nội dung khóa.
    \item \textbf{Cơ chế xoay vòng Refresh Token (Token Rotation \& Token Family):} Khi Access Token hết hạn, Front-end tự động gửi request đến \texttt{POST /api/v1.0/auth/refresh}. Backend xác thực cookie, thu hồi (revoke) Refresh Token cũ và đồng thời cấp một cặp Access Token + Refresh Token hoàn toàn mới. Nếu hệ thống phát hiện một Refresh Token đã từng bị thu hồi lại được gửi lên máy chủ (dấu hiệu của kẻ tấn công cố tình phát lại token đánh cắp được - Replay Attack), cơ chế bảo vệ sẽ lập tức vô hiệu hóa toàn bộ các token thuộc cùng nhóm phiên (\texttt{familyId}), buộc toàn bộ các phiên liên quan phải đăng nhập lại từ đầu.
\end{enumerate}

% ==============================================================================
\section{Hệ quản trị cơ sở dữ liệu và quản lý phiên bản}
\label{sec:he_quan_tri_csdl}

\subsection{Hệ quản trị cơ sở dữ liệu MySQL Server (InnoDB Engine)}
Cơ sở dữ liệu của hệ thống POS đòi hỏi tính chính xác, tính nhất quán và khả năng chịu tải giao dịch cao. Dự án lựa chọn \textbf{MySQL Server 8.0} kết hợp cùng công cụ lưu trữ mặc định \textbf{InnoDB} \cite{mysqldoc2024}. InnoDB đảm bảo đầy đủ 4 tính chất giao dịch ACID cốt lõi:
\begin{itemize}[leftmargin=1.5cm]
    \item \textbf{Nguyên tử (Atomicity):} Quá trình thanh toán đơn hàng bao gồm nhiều bước: ghi hóa đơn, ghi chi tiết đơn hàng, trừ số dư kho sổ cái, cập nhật doanh số khách hàng CRM và tăng số lượt sử dụng voucher khuyến mãi. Tất cả các bước này đều được bọc trong một khối giao dịch duy nhất (\texttt{@Transactional}). Nếu xảy ra bất kỳ lỗi ngoại lệ nào trong quá trình thực hiện, toàn bộ các thay đổi đều tự động phục hồi về trạng thái ban đầu (Rollback), không để lại dữ liệu rác.
    \item \textbf{Nhất quán (Consistency):} Ràng buộc toàn vẹn dữ liệu được thiết lập chặt chẽ bằng các khóa ngoại (Foreign Keys), ngăn chặn tình trạng xóa danh mục khi vẫn còn sản phẩm liên kết (\texttt{ON DELETE RESTRICT}), hoặc tự động xóa sạch các biến thể khi sản phẩm cha bị gỡ bỏ (\texttt{ON DELETE CASCADE}).
    \item \textbf{Cô lập (Isolation):} Mức độ cô lập giao dịch chuẩn \texttt{READ\_COMMITTED} kết hợp cơ chế khóa mức bản ghi (Row-level Locking) của InnoDB giúp nhiều nhân viên thu ngân có thể thực hiện thanh toán song song mà không gây ra xung đột dữ liệu tồn kho.
    \item \textbf{Bền vững (Durability):} Mọi giao dịch sau khi xác nhận (Commit) thành công đều được ghi nhận vào nhật ký ghi trước (Redo Log), đảm bảo dữ liệu không bị thất thoát ngay cả khi máy chủ gặp sự cố mất điện đột ngột.
\end{itemize}

\subsection{Quản lý phiên bản cấu trúc CSDL bằng Flyway Migration}
Trong quá trình phát triển phần mềm theo nhóm, việc cấu hình tự động cập nhật schema của Hibernate (\texttt{spring.jpa.hibernate.ddl-auto=update}) tiềm ẩn rủi ro rất lớn: Hibernate không thể theo dõi lịch sử thay đổi cấu trúc bảng, dễ gây ra lỗi khóa bảng hoặc thậm chí làm mất dữ liệu khi đổi tên trường trên môi trường thực tế (Production).

Để giải quyết triệt để vấn đề này, dự án tích hợp công cụ mã nguồn mở \textbf{Flyway Migration}:
\begin{itemize}[leftmargin=1.5cm]
    \item Cấu hình JPA chuyển sang chế độ kiểm tra nghiêm ngặt: \texttt{spring.jpa.hibernate.ddl-auto=validate}.
    \item Mọi thay đổi về cấu trúc bảng (thêm cột, sửa kiểu dữ liệu, đánh chỉ mục index) đều được văn bản hóa thành các tệp kịch bản SQL có đánh số phiên bản trong thư mục \texttt{src/main/resources/db/migration/} (ví dụ: \texttt{V2\_\_.sql}).
    \item Khi ứng dụng khởi động, Flyway tự động đối chiếu bảng quản lý \texttt{flyway\_schema\_history} trên MySQL để xác định các tệp kịch bản mới chưa được thực thi và tự động áp dụng tuần tự, đảm bảo tính đồng nhất tuyệt đối về cấu trúc cơ sở dữ liệu trên môi trường của tất cả các thành viên trong nhóm cũng như trên máy chủ triển khai.
\end{itemize}

% ==============================================================================
\section{Cổng thanh toán và dịch vụ đám mây tích hợp}
\label{sec:cong_thanh_toan_cloud}

\subsection{Cổng thanh toán trực tuyến PayOS (Chuẩn VietQR động)}
Tại thị trường bán lẻ Việt Nam hiện nay, thói quen thanh toán không dùng tiền mặt thông qua chuyển khoản ngân hàng quét mã QR (VietQR theo chuẩn NAPAS) đang tăng trưởng với tốc độ bùng nổ. Tuy nhiên, việc chuyển khoản thủ công truyền thống (khách hàng tự nhập số tài khoản, số tiền và nội dung chuyển khoản) thường xuyên gây ra ùn tắc tại quầy thu ngân và tạo kẽ hở cho các hành vi gian lận sửa ảnh chụp màn hình chuyển khoản.

Dự án đã tích hợp thành công giải pháp cổng thanh toán mở **PayOS** thông qua bộ thư viện chính thức \texttt{vn.payos:payos-java:2.0.1} \cite{payosdoc2024}:
\begin{itemize}[leftmargin=1.5cm]
    \item \textbf{Sinh mã VietQR động theo đơn hàng:} Khi thu ngân chọn phương thức thanh toán \texttt{PAYOS}, hệ thống tự động gọi API của PayOS để tạo một liên kết thanh toán (Payment Link) có mã định danh gắn liền với \texttt{orderCode} của đơn hàng. Mã QR được trả về và hiển thị tức thời lên màn hình thanh toán kèm chính xác số tiền cần trả và mã đơn hàng đã được khóa cứng.
    \item \textbf{Cơ chế Webhook xác thực chữ ký điện tử (HMAC-SHA256):} Khi khách hàng quét mã và tiền đã vào tài khoản ngân hàng của cửa hàng, máy chủ PayOS sẽ gửi ngay một tín hiệu HTTP POST (Webhook) đến endpoint \texttt{/api/v1.0/payos/webhook} của hệ thống. Để ngăn chặn các cuộc tấn công làm giả Webhook, backend sử dụng mã khóa bí mật \texttt{checksumKey} để thẩm định lại chữ ký điện tử đính kèm trong payload. Chỉ khi chữ ký hoàn toàn hợp lệ, đơn hàng mới được tự động chuyển sang trạng thái hoàn tất (\texttt{COMPLETED}).
    \item \textbf{Quy trình chuyển đổi thanh toán linh hoạt tại quầy (Switch to Cash - ADR-0006):} Trong trường hợp khách hàng gặp sự cố về ứng dụng ngân hàng hoặc nghẽn mạng 4G, thu ngân có thể lập tức nhấn nút \textit{''Chuyển sang tiền mặt''} trên giao diện POS. Hệ thống sẽ kích hoạt API \texttt{POST /orders/\{orderId\}/switch-to-cash}, chuyển trạng thái đơn hàng sang \texttt{COMPLETED} với phương thức tiền mặt và tự động gửi lệnh hủy liên kết thanh toán từ xa trên cổng PayOS mà không làm nhân đôi xuất kho hay tạo đơn trùng lặp.
\end{itemize}

\subsection{Dịch vụ lưu trữ đám mây Amazon S3 (AWS Simple Storage Service)}
Trong các ứng dụng Web truyền thống, hình ảnh sản phẩm thường được lưu trữ trực tiếp trên ổ đĩa cục bộ của máy chủ Web (Local Disk Storage). Cách tiếp cận này bộc lộ nhiều nhược điểm: làm tăng dung lượng sao lưu của mã nguồn, gây khó khăn khi triển khai ứng dụng bằng các Docker container (vốn có tính chất vô trạng thái - Stateless), và làm tiêu hao băng thông quý giá của ứng dụng vào việc phục vụ các tệp hình ảnh tĩnh dung lượng lớn.

Hệ thống POS giải quyết triệt để bài toán này bằng cách tích hợp dịch vụ lưu trữ đối tượng đám mây \textbf{AWS S3} thông qua bộ công cụ phát triển phần mềm chính thức **AWS SDK for Java v2** (\texttt{software.amazon.awssdk:s3:2.30.31}) \cite{awss3doc2024}:
\begin{itemize}[leftmargin=1.5cm]
    \item Khi người quản trị tải ảnh sản phẩm hoặc danh mục lên từ giao diện quản lý, dịch vụ \texttt{FileUploadServiceImpl} tiếp nhận tệp đa phần (MultipartFile), chuẩn hóa định dạng, sinh mã định danh duy nhất (UUID) cho tệp và tải trực tiếp lên kho chứa S3 Bucket thông qua thao tác \texttt{PutObjectRequest} với quyền truy cập \texttt{public-read}.
    \item Cơ sở dữ liệu MySQL chỉ cần lưu trữ chuỗi đường dẫn URL trực tiếp tới tệp trên hạ tầng của AWS (ví dụ: \texttt{https://my-bucket.s3.ap-southeast-1.amazonaws.com/uuid.jpg}). Trình duyệt của khách hàng và thu ngân sẽ tải trực tiếp hình ảnh từ hạ tầng máy chủ phân tán (CDN) của Amazon, giải phóng hoàn toàn tải trọng truyền tải tĩnh cho máy chủ Spring Boot.
    \item Dịch vụ cũng cài đặt sẵn cơ chế dọn dẹp tự động (\texttt{deleteFile}): khi một sản phẩm hoặc danh mục bị chỉnh sửa ảnh mới hoặc bị xóa khỏi hệ thống, ảnh cũ trên S3 Bucket sẽ được tự động xóa bỏ để tránh lãng phí chi phí lưu trữ đám mây.
\end{itemize}

% ==============================================================================
\section{Môi trường phát triển và triển khai mã nguồn mở}
\label{sec:moi_truong_trien_khai}

\subsection{Đóng gói ứng dụng với Docker và Docker Compose}
Để giải quyết bài toán kinh điển trong phát triển phần mềm \textit{''Chạy được trên máy của tôi nhưng lỗi trên máy khác''} do sự xung đột về phiên bản Java, Node.js hoặc cấu hình hệ điều hành cục bộ, dự án tiến hành đóng gói toàn bộ hệ thống bằng công nghệ container hóa **Docker** \cite{docker2024}:
\begin{itemize}[leftmargin=1.5cm]
    \item \textbf{Dockerfile Front-end:} Sử dụng kỹ thuật đóng gói đa giai đoạn (Multi-stage build). Giai đoạn 1 sử dụng image \texttt{node:20-alpine} để cài đặt thư viện và biên dịch mã nguồn React/Vite ra thư mục tĩnh \texttt{dist}. Giai đoạn 2 sao chép các tệp tĩnh này sang một image siêu nhẹ \texttt{nginx:alpine}, tạo ra container thành phẩm chỉ có kích thước khoảng $30$ MB, vừa bảo mật vừa tối ưu tốc độ nạp.
    \item \textbf{Dockerfile Back-end:} Sử dụng image nền tảng \texttt{eclipse-temurin:25-jdk-alpine} hỗ trợ đầy đủ môi trường thực thi Java 25, đóng gói ứng dụng Spring Boot thành một container độc lập.
    \item \textbf{Docker Compose (\texttt{docker-compose.yml}):} Định nghĩa và điều phối mạng nội bộ (Bridge Network) kết nối giữa 3 dịch vụ: \texttt{frontend}, \texttt{backend} và \texttt{mysql-db}. Toàn bộ hệ thống có thể được khởi chạy chỉ với một câu lệnh duy nhất \texttt{docker-compose up -d}.
\end{itemize}

\subsection{Máy chủ Web và Reverse Proxy Nginx}
Máy chủ **Nginx** được bố trí ở vị trí tiền tuyến (Frontline) trong kiến trúc triển khai container Front-end:
\begin{itemize}[leftmargin=1.5cm]
    \item \textbf{Phục vụ tệp tĩnh hiệu năng cao:} Xử lý việc nạp các tệp HTML, CSS, JavaScript của giao diện React với cơ chế bộ nhớ đệm (Browser Caching) và nén dữ liệu qua thuật toán Gzip, giúp giảm thiểu độ trễ mạng tại quầy thu ngân.
    \item \textbf{Xử lý định tuyến ứng dụng đơn trang (SPA Routing):} Đối với các ứng dụng Single Page Application xây dựng bằng React Router, khi người dùng tải lại trang tại một đường dẫn con (như \texttt{/explore} hay \texttt{/dashboard}), Nginx được cấu hình chỉ thị \texttt{try\_files \$uri \$uri/ /index.html} để luôn chuyển quyền điều hướng lại cho React Router xử lý, tránh lỗi 404 Not Found thường gặp.
\end{itemize}

\subsection{Hệ thống quản lý phiên bản mã nguồn Git và nền tảng GitHub}
Toàn bộ mã nguồn của dự án được quản lý phiên bản phân tán bằng công cụ **Git** và lưu trữ trên nền tảng **GitHub**:
\begin{itemize}[leftmargin=1.5cm]
    \item Nhóm áp dụng quy trình phân nhánh tính năng (Git Feature Branch Workflow): Nhánh \texttt{main} luôn duy trì trạng thái mã nguồn ổn định nhất đã được kiểm thử; mỗi thành viên khi phát triển một tính năng mới (ví dụ: tích hợp PayOS hay xây dựng trang khuyến mãi) đều tạo một nhánh riêng biệt (\texttt{feature/...}) để làm việc độc lập.
    \item Các thay đổi trước khi được gộp (Merge) vào nhánh chính đều phải thông qua quy trình đánh giá mã nguồn (Code Review), đảm bảo mã nguồn tuân thủ chặt chẽ các tiêu chuẩn mã hóa và các quyết định kiến trúc (ADR) đã được nhóm thống nhất từ đầu.
\end{itemize}

% ==============================================================================
\section{Tổng kết chương 1}
Chương 1 đã trình bày một cách toàn diện và có hệ thống về bối cảnh đề tài, sự cần thiết khách quan của việc chuyển dịch hệ thống bán lẻ POS lên nền tảng Web phân tán, cũng như cơ sở khoa học của việc lựa chọn các giải pháp công nghệ mã nguồn mở hàng đầu trong dự án. Sự kết hợp chặt chẽ giữa tính năng mượt mà, phản hồi dưới 1 giây của React 19 / Vite / Tailwind CSS tại tầng Front-end và sự ổn định, an toàn, chuẩn xác của Java 25 / Spring Boot 4.x / MySQL tại tầng Back-end, cùng các dịch vụ đám mây tiên tiến như PayOS VietQR và AWS S3, đã tạo nên một nền tảng kỹ thuật vững chắc cho việc phân tích, thiết kế và cài đặt chi tiết hệ thống trong các chương tiếp theo.
